\documentclass[11pt]{article}

\usepackage[a4paper,margin=1in]{geometry}
\usepackage[T1]{fontenc}
\usepackage[utf8]{inputenc}
\usepackage{lmodern}
\usepackage{amsmath,amssymb,amsthm}
\newtheorem{lemma}{Lemma}
\usepackage[hidelinks]{hyperref}
\usepackage{xcolor}
\usepackage{enumitem}


\setlist[itemize]{topsep=4pt,itemsep=3pt}
\setlist[enumerate]{topsep=4pt,itemsep=3pt}

\title{Next Steps for the Cut-and-Project Paper}
\author{Dirk Kunert}
\date{\today}

\begin{document}

\maketitle

\section*{Purpose of this note}

This document collects important hints, warnings, and concrete next steps
for improving the manuscript

\begin{center}
\emph{Period Length of One-Dimensional Cut-and-Project Sequences
with Rational Slope and Arbitrary Real Window Width}.
\end{center}

The central result appears to be a clean and potentially publishable formula
for the period length of a rational one-dimensional cut-and-project
construction.  However, several issues should be clarified before submitting
the paper to arXiv or to a journal.

\section*{Overall assessment}

The work is potentially worth putting on arXiv, but not yet in exactly its
current form.

\begin{itemize}
    \item The mathematical value is real, but probably modest rather than
    major.

    \item The result is more naturally described as elementary combinatorics,
    discrete geometry, or symbolic dynamics than as deep number theory.

    \item The formula is attractive because it is exact, simple, and covers
    all rational slopes and all real window widths \(\omega \ge 0\).

    \item The theorem is likely best framed as a clean explicit result about
    rational cut-and-project constructions, rather than as a major
    contribution to the general theory of aperiodic order.

    \item An arXiv preprint is reasonable after revision.

    \item A short journal note might be possible if the statement is correct,
    the proof is fully polished, and the relation to existing literature is
    handled carefully.
\end{itemize}

\section*{Central mathematical result}

The main formula is:

\[
\lambda =
\begin{cases}
N, & D \nmid N,\\
N/D, & D \mid N,
\end{cases}
\]

where

\[
N = \lfloor \omega\alpha \rfloor
  + \lfloor \omega\beta \rfloor + 1,
\qquad
D = \alpha^2+\beta^2.
\]

This is a strong feature of the paper: the result is short, exact, and easy
to state.

The proof reduces the geometry to arithmetic progressions modulo \(D\).
The key invariant is

\[
r = \alpha x - \beta y,
\]

and the projected coordinate is controlled by

\[
\tilde p = \beta x + \alpha y.
\]

For fixed \(r\), the projected values lie in an arithmetic progression modulo
\(D\).  The formula follows from counting how \(N\) consecutive values of
\(r\) distribute among the residue classes modulo \(D\).

\section*{Most important issue: multiset versus set}

The manuscript currently defines the sorted projected values as a multiset.
This is mathematically legitimate, but it must be stated extremely clearly.

In standard cut-and-project terminology, one usually speaks about point sets,
where duplicate projected points are identified.  In your manuscript,
however, different lattice points may project to the same \(\tilde p\)-value,
and these repeated values are counted with multiplicity.

This distinction is crucial.

\subsection*{Why this matters}

In the degenerate case \(D \mid N\), your proof says that every integer
\(\tilde p\) is achieved with multiplicity \(N/D\).  Hence the sorted sequence
has the form

\[
\ldots,
\underbrace{t,\ldots,t}_{N/D},
\underbrace{t+1,\ldots,t+1}_{N/D},
\ldots
\]

and the difference sequence has the form

\[
\underbrace{0,\ldots,0}_{N/D-1},1,
\underbrace{0,\ldots,0}_{N/D-1},1,\ldots
\]

This has period \(N/D\).

But this is only true if multiplicities are retained.  If duplicate projected
points are collapsed, then the set of projected values is simply all integers,
and the ordinary gap sequence has constant gap \(1\), hence period \(1\).

Therefore the theorem should explicitly be described as a theorem about the
projected lattice-point multiset, or about the multiset gap sequence.

\subsection*{Recommended fix}

Change the title, abstract, and definitions so that the word ``multiset''
appears prominently.

Possible titles:

\begin{itemize}
    \item \emph{Period Length of the Multiset Gap Sequence in Rational
    One-Dimensional Cut-and-Project Projections}

    \item \emph{Period Lengths for Projected Lattice-Point Multisets from
    Rational Cut-and-Project Strips}

    \item \emph{A Period Formula for Rational Cut-and-Project Strip
    Projections with Multiplicity}
\end{itemize}

Add a remark such as:

\begin{quote}
The projected values are counted with multiplicity throughout this paper.
Thus, if two distinct lattice points have the same projected coordinate, they
contribute two consecutive entries to the sorted projected sequence and hence
may produce a zero gap.  If multiplicities are discarded, the resulting
ordinary point-set gap sequence may have a smaller period.  In particular, in
the case \(D \mid N\), the set-valued projected sequence has constant gap and
period \(1\), whereas the multiset sequence has period \(N/D\).
\end{quote}

This clarification is essential before submission.

\section*{Tone and positioning}

The paper should be presented as a clean explicit formula in a rational
special case, not as a major breakthrough in aperiodic order.

The current introduction is acceptable in spirit, but it should be toned down
slightly.

Recommended positioning:

\begin{quote}
The irrational case is central in the theory of aperiodic order.  The
rational case is periodic and therefore often regarded as elementary or
degenerate from that perspective.  Nevertheless, the exact period length of
the projected multiset gap sequence appears not to be commonly recorded in
this explicit form.  The present note gives a complete elementary formula.
\end{quote}

Avoid statements that sound too grand, such as suggesting that the paper
fills a major gap in the theory of model sets.  The result is valuable
because it is exact and explicit, not because it revolutionises the field.

\section*{arXiv suitability}

The paper is suitable for arXiv after revision.

Recommended arXiv categories:

\begin{itemize}
    \item Primary: \texttt{math.CO} or \texttt{math.DS}
    \item Possible cross-listing: \texttt{math.NT}, \texttt{math.MG}
\end{itemize}

The work is probably not best submitted primarily as number theory.  It uses
number-theoretic tools such as congruences, residue classes, and coprimality,
but the main object is a gap sequence arising from a geometric/combinatorial
construction.

\section*{Important arXiv warning: AI co-authors}

Do not list AI systems as co-authors.

The current manuscript lists Gemini and Claude as authors.  This should be
changed before submission.

Use:

\begin{verbatim}
\author{Dirk Kunert}
\end{verbatim}

Then move the AI discussion to the acknowledgements.

Suggested wording:

\begin{quote}
The author used ChatGPT, Gemini, and Claude as AI-assisted tools for
exploration, drafting, proof checking, and Lean formalisation.  The author
has independently checked the mathematical content and accepts full
responsibility for the results and any errors.
\end{quote}

Avoid wording such as ``Gemini derived and proved the formula'' in the author
block.  That makes the paper look less professionally controlled.  The
responsible author should be the human mathematician.

\section*{Suggested revised abstract}

A safer and more professional abstract would be:

\begin{quote}
We study the gap sequence of the projected lattice-point multiset obtained
from a rational-slope strip in the plane. Let the slope be
\(\alpha/\beta\), with \(\gcd(\alpha,\beta)=1\), and let the window parameter
be \(\omega \ge 0\). Counting projected points with multiplicity, the sequence
of consecutive gaps is periodic. We prove that its minimal period is
\[
\lambda =
\begin{cases}
N, & D\nmid N,\\
N/D, & D\mid N,
\end{cases}
\]
where
\[
N=\lfloor\omega\alpha\rfloor+\lfloor\omega\beta\rfloor+1,
\qquad
D=\alpha^2+\beta^2.
\]
The proof reduces the geometry to the distribution of \(N\) consecutive
residue classes modulo \(D\). In the non-uniform case, the heavy residue
classes form a cyclic interval with trivial stabiliser, which implies
minimality of the period. We also indicate how the formula changes when
multiplicities are discarded.
\end{quote}

\section*{Proof issues to check carefully}

The proof is plausible, but the following points should be checked line by
line.

\subsection*{1. Sorting of a bi-infinite multiset}

The manuscript sorts the projected values into a bi-infinite sequence

\[
(\tilde p^{(i)}_s)_{i\in\mathbb Z}.
\]

This should be justified carefully.  One needs to know that the projected
values are locally finite as a multiset and can be indexed increasingly over
\(\mathbb Z\).

Suggested addition:

\begin{quote}
Since the projected values are a finite union of arithmetic progressions of
common difference \(D\), counted with finite multiplicities, the resulting
multiset is locally finite and unbounded in both directions.  Hence it admits
a nondecreasing enumeration indexed by \(\mathbb Z\), unique up to translation
of the index.
\end{quote}

\subsection*{2. Multiplicity convention}

State whether the sorted sequence is strictly increasing or nondecreasing.
With multiplicities, it is nondecreasing.

Use ``nondecreasing'' rather than ``ascending'' if duplicates are present.

\subsection*{3. The assertion \(\lambda\mid N\)}

The proof says that the sequence has a repeating block of length \(N\), hence
\(\lambda\mid N\).  This is correct if the block structure is precisely
periodic.  It should be made explicit that after \(N\) entries the sorted
multiset is translated by \(D\):

\[
\tilde p_s^{(i+N)} = \tilde p_s^{(i)} + D.
\]

Then

\[
\delta^{(i+N)}=\delta^{(i)}.
\]

This gives \(N\) as a period, hence the minimal period divides \(N\).

\subsection*{4. The minimality argument}

The minimality proof is the heart of the paper.  It is good, but it should be
made maximally rigorous.

The argument is:

\begin{enumerate}
    \item Suppose \(\lambda'<N\) is a period.
    \item Since \(N\) is a period, the minimal period divides \(N\), so
    \(N=k\lambda'\) for some \(k>1\).
    \item Let \(\sigma\) be the sum of one \(\lambda'\)-block of gaps.
    \item Then translation by \(\sigma\) preserves the multiset.
    \item Since \(N\) gaps sum to \(D\), we get \(k\sigma=D\).
    \item Therefore translation by \(\sigma\) induces a nontrivial
    stabiliser of the residue multiplicity function modulo \(D\).
    \item But the heavy set is a proper cyclic interval, and such an interval
    has trivial stabiliser.
    \item Contradiction.
\end{enumerate}

This is elegant.  The paper should isolate this as a lemma, perhaps:

\begin{lemma}
Let \(S\) be a finite union of arithmetic progressions of common difference
\(D\), with residue multiplicity function \(\mu\).  If the gap sequence has
period \(L\), then the translation by the sum of one \(L\)-block of gaps
preserves \(\mu\).
\end{lemma}

This would make the proof more modular.

\subsection*{5. The cyclic interval stabiliser}

The statement that a proper nonempty cyclic interval has trivial stabiliser is
correct, but the proof should be made more formal.

Suggested lemma:

\begin{lemma}
Let \(D\ge 2\), and let
\[
H=\{a,a+1,\ldots,a+s-1\}\subset \mathbb Z/D\mathbb Z
\]
with \(0<s<D\).  If \(H+\tau=H\), then \(\tau\equiv0\pmod D\).
\end{lemma}

A clean proof uses the boundary of \(H\): the set \(H\) has exactly two
boundary edges in the cyclic graph \(C_D\), whereas a nontrivial translation
moves the boundary.  Alternatively, use cardinality of intersections.

\section*{Lean formalisation}

The Lean formalisation is a strength, but the wording should be precise.

Do not overclaim that the entire geometric theorem has been formalised from
first principles if the construction is partly abstracted through a typeclass.

Recommended wording:

\begin{quote}
The residue-class and minimal-period argument has been formalised in Lean 4.
The geometric reduction from the strip construction to the residue multiset is
represented at the combinatorial level.  Thus the formalisation verifies the
algebraic core of the proof rather than the full analytic geometry from first
principles.
\end{quote}

If the repository is public, make sure that:

\begin{itemize}
    \item it builds from a clean clone;
    \item the Lean version is specified;
    \item the Mathlib commit is specified;
    \item there are no \texttt{sorry};
    \item instructions are included in a \texttt{README.md};
    \item the file names and line numbers in the paper are still correct.
\end{itemize}

\section*{Computational verification}

The computational verification is useful, but it should not be overemphasised.
Once the proof is complete, the tests are supporting evidence, not part of the
mathematical proof.

Recommended placement:

\begin{itemize}
    \item keep a short remark in the main text;
    \item move details to an appendix or repository;
    \item describe the testing method clearly enough to be reproducible.
\end{itemize}

Mention:

\begin{itemize}
    \item ranges of \(\alpha,\beta,\omega\);
    \item how \(\omega\) was sampled;
    \item how the period was computed;
    \item whether multiplicities were counted;
    \item how degenerate cases \(D\mid N\) were generated.
\end{itemize}

\section*{Literature section}

The related-literature section should be strengthened.

The current discussion of the three-distance theorem, Christoffel words, and
Sturmian sequences is useful, but the paper should avoid relying on Wikipedia
as a research citation.

Recommended changes:

\begin{itemize}
    \item Replace the Wikipedia citation for the three-distance theorem with
    a mathematical source.

    \item Keep Berstel--Lauve--Reutenauer--Saliola for Christoffel words.

    \item Keep Haynes--Koivusalo--Walton--Sadun for cut-and-project gap
    problems.

    \item Add a sentence explaining exactly why these papers do not contain
    your formula.

    \item Be cautious with the claim ``no prior work was found''.  Better:
    ``The author is not aware of a reference that records this exact period
    formula in the above multiset convention.''
\end{itemize}

Suggested wording:

\begin{quote}
The author is not aware of a reference that states the above formula for the
minimal period of the multiset gap sequence associated with a rational
one-dimensional cut-and-project strip.  The result is elementary once the
appropriate residue-class description is introduced, but it appears useful to
record explicitly.
\end{quote}

This is modest and credible.

\section*{Possible mathematical extensions}

The current open problems are good.  They can be sharpened.

\subsection*{1. Window translation}

Study what happens if the window is shifted transversally.  Does the same
formula depend only on the number \(N\) of admissible invariant values, or do
endpoint effects change the period?

\subsection*{2. Set-valued version}

Derive the exact period when projected points are counted without
multiplicity.  This may be a natural companion theorem and would remove
possible confusion with standard model-set terminology.

\subsection*{3. Gap-value distribution}

The paper determines the period length, but not the actual distribution of
gap sizes within one period.

A useful extension would be:

\begin{quote}
For fixed \(\alpha,\beta,\omega\), determine the distinct gap values and
their multiplicities in one minimal period.
\end{quote}

This would connect more directly to the three-distance theorem and gap
problems.

\subsection*{4. Higher dimensions}

This is a natural but much harder direction.  It should be described as
speculative.

\section*{Recommended structure for a revised paper}

A cleaner revised structure could be:

\begin{enumerate}
    \item Introduction
    \item Statement of the multiset convention
    \item Geometry-to-residue reduction
    \item The residue multiplicity function
    \item Generic case \(D\nmid N\)
    \item Uniform case \(D\mid N\)
    \item Set-valued comparison
    \item Computational checks
    \item Lean formalisation
    \item Related literature and open problems
\end{enumerate}

\section*{Concrete checklist before arXiv}

Before submitting to arXiv, do the following:

\begin{enumerate}
    \item Remove AI systems from the author list.

    \item Put AI assistance only in the acknowledgements.

    \item Add an explicit multiset convention near the beginning.

    \item Change the title to mention ``multiset'' or ``with multiplicity''.

    \item Replace ``ascending'' by ``nondecreasing'' where multiplicities are
    possible.

    \item Add a remark explaining the difference between the multiset gap
    sequence and the ordinary set-valued gap sequence.

    \item Make the proof of \(\tilde p_s^{(i+N)}=\tilde p_s^{(i)}+D\)
    explicit.

    \item Isolate the minimality argument as a clean lemma.

    \item State and prove the trivial-stabiliser lemma formally.

    \item Tone down claims of novelty.

    \item Replace Wikipedia citations with mathematical references.

    \item Check all bibliography entries.

    \item Check all Lean repository links.

    \item Ensure the Lean code builds from a clean clone.

    \item Make clear exactly what the Lean proof does and does not formalise.

    \item Consider adding a short section on the version without
    multiplicities.

    \item Decide the arXiv category, probably \texttt{math.CO} or
    \texttt{math.DS}, with possible cross-listing to \texttt{math.NT}.

    \item Compile the paper locally and check for overfull boxes,
    undefined references, and bibliography problems.

    \item Ask at least one human mathematician to read the proof.
\end{enumerate}

\section*{Final judgement}

The work is worth developing.

It is not likely to be regarded as a major new theorem in number theory, but
it can be a respectable and useful short mathematical note.  Its strengths are:

\begin{itemize}
    \item a clean exact formula;
    \item a transparent residue-class proof;
    \item coverage of all real \(\omega\ge0\);
    \item computational verification;
    \item partial Lean formalisation.
\end{itemize}

Its weaknesses are:

\begin{itemize}
    \item possible confusion between point sets and multisets;
    \item elementary depth;
    \item need for better literature positioning;
    \item problematic AI co-author presentation;
    \item risk of overclaiming the scope of the Lean formalisation.
\end{itemize}

With the revisions listed above, the paper is suitable for arXiv as a modest,
clear, and potentially useful contribution.

\end{document}